\renewcommand{\chaptername}{بخش}
\chapter{اتوانکدرهای واریاسیونال (VAE)}

\section{مقدمه}
اتوانکدر واریاسیونال (Variational Autoencoder یا VAE) چارچوبی برای یادگیری مدل‌های مولد با استفاده از تقریب واریانسی و شبکه‌های عصبی است. در مقابل اتوانکدرهای کلاسیک که صرفاً فشرده‌سازی-بازسازی انجام می‌دهند، VAE یک مدل احتمال‌محور بر روی فضای داده می‌سازد که امکان نمونه‌زنی صریح و تحلیل فضای پنهان را فراهم می‌نماید.

\section{فرمول‌بندی ریاضی و ELBO}
فرض کنید نمونهٔ مشاهده‌شده $x$ و متغیر پنهان $z$ داریم. مدل کامل به‌صورت توزیع مشترک $p_\\theta(x,z)=p_\\theta(x\\mid z)p(z)$ تعریف می‌شود. هدف حداکثر کردن لگ-احتمال مارژینال $\\log p_\\theta(x)$ است:
$$
\\log p_\\theta(x)=\\log\\int p_\\theta(x,z) \, dz.
$$
برای جلوگیری از محاسبهٔ مستقیمِ انتگرال، توزیع تقریبی $q_\\phi(z\\mid x)$ معرفی می‌شود و با استفاده از رابطهٔ پایه‌ٔ KL، داریم:
$$
\\log p_\\theta(x)=\\mathcal{L}(\\theta,\\phi;x)+D_{\\mathrm{KL}}(q_\\phi(z\\mid x)\\|p_\\theta(z\\mid x)),
$$
که
$$
\\mathcal{L}(\\theta,\\phi;x)=\\mathbb{E}_{q_\\phi(z\\mid x)}[\\log p_\\theta(x\\mid z)]-D_{\\mathrm{KL}}(q_\\phi(z\\mid x)\\|p(z)).
$$
از آن‌جا که KL غیرمنفی است، $\\mathcal{L}$ یک حد پایین برای لگ‌احتمال است (Evidence Lower Bound — ELBO). هدف ما بیشینه‌سازی ELBO نسبت به پارامترهای $\\theta$ و $\\phi$ است.

\subsection{تجزیهٔ ELBO و مفهوم آن}
جملهٔ اول (امیدِ لگ-احتمال شرطی) معیاری برای کیفیت بازسازی است؛ جملهٔ دوم (KL) نقشِ منظم‌سازی و فشار به تطبیق توزیع تقریبی با prior دارد. در عمل، تعادل این دو جمله تعیین‌کنندهٔ trade-off بین کیفیت بازسازی و ساختار فضای پنهان است.

\section{روش باز-پارامتری‌سازی و مشتق‌گیری}
اگر $q_\\phi(z\\mid x)$ را گاوسی با پارامترهای وابسته به $x$ در نظر بگیریم ($q_\\phi(z\\mid x)=\\mathcal{N}(\\mu_\\phi(x),\\mathrm{diag}(\\sigma^2_\\phi(x)))$)، می‌توان نمونه‌گیری را با ترفند باز-پارامتری‌سازی انجام داد:
$$
z=\\mu_\\phi(x)+\\sigma_\\phi(x)\\odot\\epsilon,\\qquad\\epsilon\\sim\\mathcal{N}(0,I).
$$
این تبدیل اجازه می‌دهد تا مشتق‌گیری از ELBO نسبت به پارامترهای $\\phi$ با استفاده از محاسبهٔ مشتق‌های نقطه‌ای انجام شود.

\section{الگوریتم آموزش}
آموزش VAE با گرادیان کاهشی روی ELBO انجام می‌شود؛ گام‌های کلی:
\begin{enumerate}
  \item برای هر مینی‌بچ از داده‌ها، رمزنگار $\\mu_\\phi(x),\\sigma_\\phi(x)$ را محاسبه کن.
  \item نمونه‌گیری $z$ با باز-پارامتری‌سازی انجام می‌شود و سپس رمزگشا $p_\\theta(x\\mid z)$ را برای بازسازی محاسبه کن.
  \item ELBO را محاسبه کن: جملهٔ بازسازی (مثلاً لگ-احتمال شرطی) و KL تحلیلی بین دو گاوس را محاسبه کن.
  \item گرادیان‌های پارامترها را محاسبه و به‌روزرسانی انجام بده (مثلاً با Adam).
\end{enumerate}

الگوریتم به صورت خلاصه:
\begin{algorithm}[H]
\caption{آموزش VAE}
\begin{enumerate}
  \item مقداردهی پارامترها $\\theta,\\phi$
  \item تا همگرایی:
  \begin{itemize}
    \item نمونه از داده $x^{(i)}$ در مینی‌بچ
    \item محاسبه $\\mu_\\phi(x),\\sigma_\\phi(x)$
    \item نمونه‌گیری $z$ با باز-پارامتری‌سازی
    \item محاسبه ضرر: $-\\mathcal{L}(\\theta,\\phi;x^{(i)})$
    \item به‌روزرسانی پارامترها با گرادیان
  \end{itemize}
\end{enumerate}
\end{algorithm}

\section{گزینه‌های طراحی و عملی}
\subsection{prior بر روی $z$}
معمولاً از prior سادهٔ استاندارد $p(z)=\\mathcal{N}(0,I)$ استفاده می‌شود. در مواردی که prior پیچیده‌تر لازم است، از جریان‌های نرمالیزاسیون (Normalizing Flows) برای افزایش انعطاف‌پذیری $q_\\phi$ استفاده می‌کنند.

\subsection{توابع بازسازی}
برای تصاویر با پیکسل‌های باینری معمولاً از برنولی و برای تصاویر پیوسته از گاوسی استفاده می‌شود. انتخاب توزیع بازسازی بر معیار لگ-احتمال تأثیر می‌گذارد.

\subsection{معماری‌های پیشنهادی}
برای تصاویر کوچک مانند dSprites از معماری‌های سبک کانولوشن (سه تا چهار لایه) و برای تصاویر پیچیده‌تر از شبکه‌های عمیق‌تر و بلوک‌های ResNet استفاده می‌شود.

\section{بسط‌ها و نسخه‌های پیشرفته}
\subsection{beta-VAE}
در beta-VAE، جملهٔ KL با ضریب $\\beta$ وزن‌دهی می‌شود:
$$
\\mathcal{L}_\\beta=\\mathbb{E}_{q} [\\log p_\\theta(x\\mid z)]-\\beta D_{\\mathrm{KL}}(q_\\phi(z\\mid x)\\|p(z)).
$$
مقدار بزرگ‌تر $\\beta$ باعث افزایش فشردگی و جداسازی عوامل می‌شود ولی می‌تواند بازسازی را تضعیف کند.

\subsection{Normalizing Flows}
جریان‌های نرمالیزاسیون با اعمال نگاشت‌های برگشت‌پذیر روی نمونهٔ اولیه، توزیع‌های پیچیده‌تری را می‌سازند و می‌توانند قدرت مدل را افزایش دهند (مثلاً IAF, RealNVP).

\section{معیارهای ارزیابی دقیق‌تر}
\subsection{MIG (Mutual Information Gap)}
MIG برای هر عامل مولد حقیقی $v$ مقدار اطلاعات متقابل $I(v; z_j)$ را محاسبه می‌کند، سپس اختلاف بین عضو اول و دوم را نرمال‌سازی می‌کند و میانگین گرفته می‌شود:
$$
\\mathrm{MIG}=\\frac{1}{|V|}\\sum_{v\\in V}\\frac{1}{H(v)}(I(v; z_{j^*})-I(v; z_{j^{**}})).
$$

\subsection{Other metrics}
DCI, FactorVAE score و سایر معیارها می‌توانند به همراه بررسی‌های بصری استفاده شوند.

\section{نکات عملی برای آموزش}
\begin{itemize}
  \item نرمال‌سازی داده‌ها و استفاده از batch normalization/LayerNorm در صورت لزوم.
  \item مقداردهی اولیه مناسب و نرخ یادگیری معقول (1e-3 یا 5e-4 برای Adam).
  \item پایش تعداد epochها و استفاده از early-stopping بر اساس معیار بازسازی یا اعتبارسنجی.
  \item ثبت مقادیر KL و reconstruction به‌صورت جداگانه برای تحلیل trade-off.
\end{itemize}

\section{نمونه کد مختصر (نمایشی)}
\begin{program}[H]
\begin{minted}[frame=none]{python}
# نمونهٔ بسیار خلاصه‌شده: محاسبهٔ ELBO برای یک مینی‌بچ
def elbo(model, x):
    mu, logvar = model.encoder(x)
    z = mu + (0.5*logvar).exp()*torch.randn_like(mu)
    recon = model.decoder(z)
    recon_loss = -reconstruction_log_prob(x, recon)
    kl = 0.5 * (mu.pow(2) + logvar.exp() - logvar - 1).sum(dim=1).mean()
    return (recon_loss + kl).mean()
\end{minted}
\caption{نمونهٔ سادهٔ محاسبهٔ ELBO}
\end{program}

\section{خلاصه و نتیجه‌گیری بخش}
VAEها چارچوبی انعطاف‌پذیر برای یادگیری نمایه‌های مفهومی و مدل‌های مولد ارائه می‌دهند. با تنظیمِ مناسب prior، اندازهٔ فضای پنهان و پارامترهای آموزشی می‌توان بین کیفیت بازسازی و تفکیک‌پذیری نمایه‌ها تعادل برقرار کرد. روش‌های پیشرفته مانند جریان‌های نرمالیزاسیون و تنظیمات \beta می‌توانند عملکرد را بهبود بخشند.
\renewcommand{\chaptername}{بخش}
\chapter{اتوانکدرهای واریاسیونال (VAE)}

\section{مقدمه}
اتوانکدر واریاسیونال (Variational Autoencoder یا VAE) یک چارچوب یادگیری غیرنظارت‌شده برای مدل‌سازی توزیعِ داده‌ها با استفاده از مدل‌های احتمالاتی و شبکه‌های عصبی است. VAE به کمک تقریب واریانسی یک توزیع پنهان $p(z\mid x)$ را مدل می‌کند و امکان نمونه‌زنی و تولید داده‌های جدید را فراهم می‌سازد.

\section{فرموله‌سازی و تابع هدف}
فرض کنید دادهٔ مشاهده‌شده $x$ و متغیر پنهان $z$ را داشته باشیم. توزیع مدل‌سازی‌شده از طریق مدل مولد $p_\theta(x\mid z)$ و پیش‌فرضِ گَردآور $p(z)$ تعریف می‌شود. هدف حداکثر کردن لگ-حتمیّت مارژینال $\log p_\theta(x)$ است که محاسبهٔ مستقیم آن معمولاً غیرممکن است. با معرفی توزیع تقریبی $q_\phi(z\mid x)$ (اپروکساتور واریانسی)، آنگاه رابطهٔ ELBO به‌دست می‌آید:
$$
\log p_\theta(x) \geq \mathcal{L}(\theta,\phi;x) = \mathbb{E}_{q_\phi(z\mid x)}[\log p_\theta(x\mid z)] - D_{\mathrm{KL}}\big(q_\phi(z\mid x)\ \|\ p(z)\big).
$$
بخش اول نمایانگر دقت بازسازی (reconstruction) و بخش دوم یک ضریب منظم‌کنندهٔ کل-اطلاعات است.

\section{تکنیک باز-پارامتری‌سازی}
برای امکان مشتق‌گیری و آموزش شبکه‌ها، از ترفند باز-پارامتری‌سازی استفاده می‌شود. برای مثال اگر $q_\phi(z\mid x)=\mathcal{N}(\mu_\phi(x),\sigma^2_\phi(x)I)$، نمونه‌گیری $z$ را می‌توان به صورت
$$
z = \mu_\phi(x) + \sigma_\phi(x) \odot \epsilon,\qquad \epsilon\sim\mathcal{N}(0,I)
$$
نوشت تا گرادیان‌ها بتوانند از طریق پارامترهای $\phi$ جریان یابند.

\section{VAE و نسخهٔ $\beta$-VAE}
برای کنترل بهتر فشردگی و جداپذیری نمایه‌های پنهان، تابع هزینه را می‌توان با ضریب $\beta$ تعدیل کرد:
$$
\mathcal{L}_\beta = \mathbb{E}_{q_\phi(z\mid x)}[\log p_\theta(x\mid z)] - \beta\, D_{\mathrm{KL}}(q_\phi(z\mid x)\|p(z)).
$$
مقادیر بزرگتر $\beta$ معمولاً به نمایه‌های پنهان منظم‌تر و تفکیک‌پذیرتر (disentangled) منجر می‌شوند، اما ممکن است کیفیت بازسازی کاهش یابد.

\section{معماری و انتخاب‌های عملیاتی}
معمولا VAEها دارای دو شبکهٔ مجزا هستند: رمزنگار (encoder) که $q_\phi(z\mid x)$ را برآورد می‌کند و رمزگشا (decoder) که $p_\theta(x\mid z)$ را تولید می‌کند. انتخاب ساختار شبکه (کانولوشنی برای تصاویر، یا لایه‌های تمام‌متصل برای بردارها)، اندازهٔ فضای پنهان، و تابع بازسازی (مثلا گاوسی یا چندکلاسه) از تصمیمات مهم پیاده‌سازی هستند.

\section{داده‌ها و جزئیات تجربی}
در این پروژه از مجموعهٔ دادهٔ dSprites برای ارزیابی استفاده شده است که شامل تصاویر سادهٔ دو‌بعدی با عوامل مرکزی مانند شکل، مقیاس و زاویه است. پیش‌پردازش استاندارد شامل نرمال‌سازی پیکسلی و دسته‌بندی‌های کوچک برای آموزش است. تنظیمات رایج برای آموزش:
\begin{itemize}
    \item بهینه‌ساز: Adam با نرخ یادگیری 1e-3
    \item اندازهٔ فضای پنهان: معمولاً در بازهٔ 10--32
    \item تعداد اپوک‌ها: 100--200 بسته به اندازهٔ مجموعه
\end{itemize}

\section{شاخص‌های ارزیابی}
برای ارزیابی از معیارهای زیر استفاده می‌شود:
\begin{itemize}
    \item کیفیت بازسازی: معیارهایی نظیر MSE یا نابرابری لگ-احتمال میانگین
    \item جداسازی نمایه‌ها: معیارهایی مانند MIG یا بررسی بصری تراورس‌های فضای پنهان
    \item تحلیل‌های بصری: بازسازی‌ها، نمونه‌های تولیدی و شبکه‌های تراورس
\end{itemize}

\section{نتایج نمونه‌ای و تفسیر}
با افزایش مقدار $\beta$ مشاهده می‌شود که نمایه‌های پنهان تمایل به تفکیک بهتر عوامل مولد دارند؛ اما کیفیت بازسازی کاهش می‌یابد. نمودارهای آموزشی شامل کاهش ضرر کلی و روند KL و بازسازی باید گزارش شده و تحلیل شوند. همچنین تراورس‌های فضای پنهان نشان می‌دهد که هر بعد پنهان تا چه اندازه با یک عامل مولد متناظر است.

\section{جمع‌بندی و چشم‌انداز}
VAEها چارچوبی محکم برای مدل‌سازی مولد ارائه می‌دهند که با ترکیب تقریب واریانسی و شبکه‌های عصبی امکان آموزش و نمونه‌زنی را فراهم می‌کنند. پژوهش‌های آینده می‌تواند شامل ترکیب VAEها با جریان‌های نرمالیزاسیون (Normalizing Flows)، استفاده از معیارهای قوی‌تر جداسازی و اعمال در داده‌های پیچیده‌تر باشد.
% \setcounter{chapter}{\Session-2}
\renewcommand{\chaptername}{سوال}
\chapter{عنوان/سوال/دو}
گزارش سوال دو.